% !TeX root = ../main.tex

% ==================================================
% A.11. 클래스·메서드 바인딩과 데코레이터
% ==================================================

\subsection{클래스·메서드 바인딩과 데코레이터}

Method Binding · Descriptor Protocol · classmethod · staticmethod · property · Decorator Stacking

\subsubsection{이전 학습과의 연결}

\zsh{A.10}에서는 여러 데코레이터를 다음과 같은 함수 객체 변환 과정으로 분석했다.

\begin{dhkcode}
@first
@second
def calculate(x):
    return x * 2
\end{dhkcode}

개념적으로 다음과 대응한다.

\begin{dhkcode}
calculate = first(second(calculate))
\end{dhkcode}

이 과정에서 원본 함수 객체가 생성되고, 데코레이터가 반환한 객체가 최종적으로 \zsh{calculate}라는 이름에 바인딩된다.

그런데 함수가 클래스 내부에 정의되면 새로운 문제가 등장한다.

\begin{dhkcode}
class Person:

    def introduce(self):
        return "Hello"
\end{dhkcode}

이때 \zsh{introduce}는 어떤 객체일까?

일반적인 함수 객체일까, 아니면 처음부터 인스턴스 메서드 객체일까?

그리고 다음 두 표현식은 왜 다르게 동작할까?

\begin{dhkcode}
Person.introduce
Person().introduce
\end{dhkcode}

이 차이를 설명하려면 데코레이터뿐 아니라 속성 탐색(Attribute Lookup)과 디스크립터 프로토콜(Descriptor Protocol)을 함께 이해해야 한다.

이번 학습의 핵심은 다음과 같다.

\begin{dhksmry}[centered]
데코레이터는 클래스 네임스페이스에 저장될 객체를 결정하고,

디스크립터는 객체가 속성 조회 과정에서 어떻게 바인딩되는지 결정한다.
\end{dhksmry}

두 과정은 서로 다른 시점에 발생한다.

\subsubsection{클래스 내부의 \zsh{def}는 무엇을 생성할까?}

먼저 데코레이터가 없는 클래스를 살펴보자.

\begin{dhkcode}
class Person:

    def __init__(self, name):
        self.name = name

    def introduce(self):
        return f"Hello, {self.name}"

p = Person("Kim")
\end{dhkcode}

클래스 정의문이 실행될 때 클래스 본문도 실행된다.

이 과정에서 \zsh{def introduce} 문장은 일반적인 Python 함수 객체를 생성하고, \zsh{introduce}라는 이름을 클래스 생성에 사용되는 네임스페이스에 바인딩한다.

클래스 객체가 생성된 뒤에는 다음과 같이 확인할 수 있다.

\begin{dhkcode}
print(type(Person.__dict__["introduce"]))

print(
    Person.__dict__["introduce"]
    is Person.introduce
)
\end{dhkcode}

예상 출력:

\begin{dhkzsh}
<class 'function'>
True
\end{dhkzsh}

여기서 첫 번째 결과가 중요하다.

\zsh{Person} 클래스의 \zsh{__dict__}에 저장된 것은 바운드 메서드가 아니라 함수 객체(Function Object)다.

\zsh{Person.__dict__}는 클래스 네임스페이스를 보여주는 \zsh{mappingproxy} 객체이고, 그 안에서 \zsh{"introduce"}라는 키가 함수 객체를 참조하고 있다.

반면 \zsh{p.__dict__}에는 다음과 같은 인스턴스 상태가 저장된다.

\begin{dhkcode}
print(p.__dict__)
\end{dhkcode}

예상 출력:

\begin{dhkzsh}
{'name': 'Kim'}
\end{dhkzsh}

즉, 클래스의 함수와 인스턴스의 데이터는 서로 다른 네임스페이스에 존재한다.

\begin{dhkdir}[in=6]
Person 클래스 객체
> __dict__
>> "__init__" -> 함수 객체
>> "introduce" -> 함수 객체
\end{dhkdir}

\begin{dhkdir}[in=6]
p 인스턴스
> __class__ -> Person
> __dict__
>> "name" -> "Kim"
\end{dhkdir}

그런데 \zsh{p.introduce}를 조회하면 단순히 클래스에 저장된 함수 객체를 그대로 반환하지 않는다.

왜 그럴까?

일반적인 사용자 정의 함수 객체가 \zsh{__get__}을 제공하는 비데이터 디스크립터(Non-data Descriptor)이기 때문이다.

\subsubsection{함수 객체가 바운드 메서드가 되는 과정}

다음 코드를 살펴보자.

\begin{dhkcode}
class Person:

    def __init__(self, name):
        self.name = name

    def introduce(self):
        return f"Hello, {self.name}"


p = Person("Kim")

func = Person.__dict__["introduce"]
method = p.introduce

print(type(func))
print(type(method))

print(method.__self__ is p)
print(method.__func__ is func)

print(method())
\end{dhkcode}

예상 출력:

\begin{dhkzsh}
<class 'function'>
<class 'method'>
True
True
Hello, Kim
\end{dhkzsh}

이 결과를 순서대로 분석해 보자.

\begin{dhknote}
\begin{dhksmry}[centered]
1. 클래스 네임스페이스에서 함수 객체 탐색
\end{dhksmry}
\end{dhknote}

\begin{dhkcode}
func = Person.__dict__["introduce"]
\end{dhkcode}

\zsh{func}는 원본 함수 객체를 참조한다.

\begin{dhknote}
\begin{dhksmry}[centered]
2. 인스턴스에서 속성 조회
\end{dhksmry}
\end{dhknote}

\begin{dhkcode}
method = p.introduce
\end{dhkcode}

일반적인 인스턴스 속성 조회에서는 클래스의 MRO를 따라 속성을 탐색하고, 발견된 객체가 디스크립터라면 관련 규칙에 따라 \zsh{__get__}을 호출할 수 있다.

여기서 함수 객체의 \zsh{__get__}이 사용된다.

핵심 동작은 다음과 같이 표현할 수 있다.

\begin{dhkcode}
method = func.__get__(p, Person)
\end{dhkcode}

그 결과 바운드 메서드 객체가 반환된다.

\begin{dhknote}
\begin{dhksmry}[centered]
3. 바운드 메서드의 내부 참조
\end{dhksmry}
\end{dhknote}

바운드 메서드는 다음 두 객체를 참조한다.

\begin{dhktable}
|속성|	참조 대상|
|:---|:---|
|method.__self__|	바인딩된 p 인스턴스|
|method.__func__|	원본 introduce 함수 객체|
\end{dhktable}

따라서 다음과 같은 구조가 형성된다.

\begin{dhkgraph}[uniform]
node n1 (1,2) :: p.introduce // 바운드 메서드 객체
node n2 (2,1) :: p 인스턴스 // name = "Kim"
node n3 (2,3) :: introduce 함수 // 원본 코드 객체 사용

split n1 -> n2 :: __self__,n3 :: __func__

caption :: 바운드 메서드는 바인딩된 인스턴스와 함수 객체를 연결하는 별도의 객체다.
\end{dhkgraph}

\begin{dhknote}
\begin{dhksmry}[centered]
4. 바운드 메서드 호출
\end{dhksmry}
\end{dhknote}

\begin{dhkcode}
method()
\end{dhkcode}

바운드 메서드는 자신이 참조하는 인스턴스를 원본 함수의 첫 번째 인수로 전달한다.

따라서 일반적인 경우 다음 두 호출은 같은 함수 실행을 수행한다.

\begin{dhkcode}
p.introduce()
Person.introduce(p)
\end{dhkcode}

여기서 \zsh{self}는 특별한 키워드가 아니다.

함수 정의에서 사용한 첫 번째 매개변수 이름일 뿐이다.

인스턴스를 자동 전달하는 기능은 \zsh{self}라는 이름에서 발생하는 것이 아니라, 함수 디스크립터가 만들어 낸 바운드 메서드의 호출 규칙에서 발생한다.

\subsubsection{세 종류의 메서드는 무엇이 다른가?}

다음 클래스를 보자.

\begin{dhkcode}
class Example:

    def instance_method(self):
        return self

    @classmethod
    def class_method(cls):
        return cls

    @staticmethod
    def static_method():
        return "static"


obj = Example()
\end{dhkcode}

세 메서드 모두 \zsh{def} 문장을 통해 함수 객체를 생성하지만, 데코레이터 적용 이후 클래스 네임스페이스에 저장되는 객체가 달라진다.

\begin{dhkcode}
print(type(Example.__dict__["instance_method"]))
print(type(Example.__dict__["class_method"]))
print(type(Example.__dict__["static_method"]))
\end{dhkcode}

예상 출력:

\begin{dhkzsh}
<class 'function'>
<class 'classmethod'>
<class 'staticmethod'>
\end{dhkzsh}

그리고 클래스나 인스턴스를 통해 속성을 조회하면 다음과 같은 결과를 얻는다.

\begin{dhkcode}
print(type(Example.instance_method))
print(type(obj.instance_method))

print(type(Example.class_method))
print(type(obj.class_method))

print(type(Example.static_method))
print(type(obj.static_method))
\end{dhkcode}

예상 출력:

\begin{dhkzsh}
<class 'function'>
<class 'method'>
<class 'method'>
<class 'method'>
<class 'function'>
<class 'function'>
\end{dhkzsh}

이 차이를 표로 정리하면 다음과 같다.

\begin{dhktable}
|구분|	클래스에 저장된 객체|	클래스 접근 결과|	인스턴스 접근 결과|
|:---|:---|:---|:---|
|일반 인스턴스 메서드|	함수 객체|	함수 객체|	인스턴스 바운드 메서드|
|@classmethod|	classmethod 디스크립터|	클래스 바운드 메서드	|클래스 바운드 메서드|
|@staticmethod|	staticmethod 디스크립터|	함수 객체	|함수 객체|
\end{dhktable}

이 표가 이번 학습 전체를 관통하는 기준이다.

\zsh{classmethod}와 \zsh{staticmethod}의 차이를 단순히 \zsh{self}와 \zsh{cls}의 차이로 이해하면 불완전하다.

실제 차이는 클래스 네임스페이스에 저장된 디스크립터가 속성 조회 과정에서 어떤 객체를 반환하느냐에 있다.

\subsubsection{\zsh{classmethod} - 함수 객체를 클래스에 바인딩하기}

먼저 클래스 메서드부터 분석해 보자.

\begin{dhkcode}
class Animal:

    @classmethod
    def identify(cls):
        return cls.__name__


a = Animal()
\end{dhkcode}

이 코드에서 다음 부분은

\begin{dhkcode}
@classmethod
def identify(cls):
    return cls.__name__
\end{dhkcode}

개념적으로 다음 변환식과 대응한다.

\begin{dhkcode}
def identify(cls):
    return cls.__name__

identify = classmethod(identify)
\end{dhkcode}

\zsh{classmethod}는 원본 함수 객체를 받아 디스크립터 객체를 생성한다.

따라서 클래스 네임스페이스에는 일반 함수 객체가 아니라 \zsh{classmethod} 객체가 저장된다.

\begin{dhkcode}
raw = Animal.__dict__["identify"]

print(type(raw))
print(type(raw.__func__))
\end{dhkcode}

예상 출력:

\begin{dhkzsh}
<class 'classmethod'>
<class 'function'>
\end{dhkzsh}

여기서 \zsh{raw.__func__}는 \zsh{classmethod}가 감싸고 있는 원본 함수 객체를 참조한다.

이제 클래스를 통해 속성을 조회해 보자.

\begin{dhkcode}
method = Animal.identify

print(type(method))
print(method.__self__ is Animal)
print(method.__func__ is raw.__func__)
print(method())
\end{dhkcode}

예상 출력:

\begin{dhkzsh}
<class 'method'>
True
True
Animal
\end{dhkzsh}

왜 이런 결과가 발생할까?

\zsh{Animal.identify}를 조회하면 클래스 네임스페이스에서 \zsh{classmethod} 디스크립터를 찾게 된다.

그 디스크립터의 \zsh{__get__}이 호출되면서 원본 함수 객체를 \zsh{Animal} 클래스에 바인딩한 메서드가 반환된다.

따라서 다음과 같은 관계가 형성된다.

\begin{dhkdiagram}[in=4]
node n1 (1,2) :: Animal 클래스 객체
node n2 (3,2) :: __dict__["identify"]
node n3 (5,2) :: classmethod 객체
node n4 (6,3) :: __func__
node n5 (6,4) :: 원본 identify 함수
node n6 (7,2) :: Animal.identify
node n7 (9,2) :: 바운드 메서드
node n10 (9,3) :: __func__
node n8 (10,3) :: __self__
node n9 (10,4) :: Animal 클래스 객체

edge n1 -- n2
edge n2 -> n3
edge n3 -- n4 :: VH
edge n4 -> n5
edge n6 -> n7
edge n7 -- n8 :: VH
edge n8 -> n9
edge n7 -- n10
edge n10 -> n5 :: HV
\end{dhkdiagram}

이제 다음 호출을 생각해 보자.

\begin{dhkcode}
Animal.identify()
\end{dhkcode}

바운드 메서드가 \zsh{Animal}을 첫 번째 인수로 전달하므로 원본 함수에서는 다음 바인딩이 형성된다.

\begin{dhksmry}[centered]
cls \ra Animal 클래스 객체
\end{dhksmry}

즉, 개념적으로 다음과 같은 호출이다.

\begin{dhkcode}
raw.__func__(Animal)
\end{dhkcode}

\begin{dhknote}
\begin{dhksmry}[centered]
인스턴스를 통해 조회한다면?
\end{dhksmry}
\end{dhknote}

\begin{dhkcode}
method = a.identify

print(method.__self__ is Animal)
print(method.__self__ is a)
\end{dhkcode}

예상 출력:

\begin{dhkzsh}
True
False
\end{dhkzsh}

인스턴스를 통해 접근하더라도 \zsh{classmethod}는 인스턴스가 아니라 클래스에 바인딩된다.

따라서 클래스 메서드의 핵심은 첫 번째 매개변수 이름이 \zsh{cls}라는 사실이 아니라, 디스크립터가 클래스 객체를 바인딩한다는 사실이다.

\subsubsection{상속과 \zsh{classmethod} - MRO와 동적 클래스 바인딩}

이제 이전에 학습한 MRO와 연결해 보자.

\begin{dhkcode}
class Parent:

    @classmethod
    def identify(cls):
        return cls.__name__


class Child(Parent):
    pass


print(Parent.identify())
print(Child.identify())

print(Parent.identify.__self__ is Parent)
print(Child.identify.__self__ is Child)

print(
    Parent.identify.__func__
    is Child.identify.__func__
)
\end{dhkcode}

예상 출력:

\begin{dhkzsh}
Parent
Child
True
True
True
\end{dhkzsh}

여기서 중요한 질문이 있다.

\zsh{Child}는 \zsh{identify}를 정의하지 않았는데 왜 \zsh{cls}가 \zsh{Parent}가 아니라 \zsh{Child}를 참조할까?

속성 탐색과 메서드 바인딩을 구별하면 이해할 수 있다.

먼저 \zsh{Child.identify}라는 속성을 조회할 때 MRO를 따라 

\zsh{Parent}에서 \zsh{identify} 디스크립터를 찾는다.

하지만 그 디스크립터가 반환하는 메서드는 실제 접근한 클래스인 \zsh{Child}에 바인딩된다.

따라서 다음 관계가 형성된다.

\begin{dhkdir}[in=6]
Child.identify
|바운드 메서드
> __self__ -> Child
> __func -> Parent에 정의된 identify 함수
\end{dhkdir}

즉, 함수를 어디에서 찾았는지와 그 함수가 어느 클래스에 바인딩되는지는 별개의 문제다.

이것이 클래스 메서드를 대체 생성자(Alternative Constructor)로 사용할 수 있는 중요한 이유다.

\begin{dhkcode}
class Model:

    def __init__(self, hidden_size):
        self.hidden_size = hidden_size

    @classmethod
    def from_default(cls):
        return cls(128)


class TransformerModel(Model):
    pass


model = TransformerModel.from_default()

print(type(model).__name__)
print(model.hidden_size)
\end{dhkcode}

예상 출력:

\begin{dhkzsh}
TransformerModel
128
\end{dhkzsh}

\zsh{from_default}는 \zsh{Model}에서 정의되었지만

\zsh{TransformerModel.from_default()}로 호출했으므로 \zsh{cls}는

\zsh{TransformerModel}을 참조한다.

따라서 다음 표현식은

\begin{dhkcode}
cls(128)
\end{dhkcode}

\zsh{TransformerModel(128)}을 호출하는 효과를 가진다.

다만 하위 클래스가 생성자 인수 구조를 변경했다면 \zsh{cls(128)}이 반드시 성공하는 것은 아니다.

\subsubsection{\zsh{staticmethod} - 자동 바인딩을 수행하지 않는 디스크립터}

이번에는 정적 메서드를 살펴보자.

\begin{dhkcode}
class Calculator:

    @staticmethod
    def add(a, b):
        return a + b


obj = Calculator()
\end{dhkcode}

클래스 본문의 다음 부분은

\begin{dhkcode}
@staticmethod
def add(a, b):
    return a + b
\end{dhkcode}

개념적으로 다음과 대응한다.

\begin{dhkcode}
def add(a, b):
    return a + b

add = staticmethod(add)
\end{dhkcode}

클래스 네임스페이스에는 \zsh{staticmethod} 객체가 저장된다.

\begin{dhkcode}
raw = Calculator.__dict__["add"]

print(type(raw))
print(type(raw.__func__))
\end{dhkcode}

예상 출력:

\begin{dhkzsh}
<class 'staticmethod'>
<class 'function'>
\end{dhkzsh}

하지만 클래스나 인스턴스를 통해 속성을 조회하면 감싸고 있던 함수 객체가 반환된다.

\begin{dhkcode}
print(Calculator.add is raw.__func__)
print(obj.add is raw.__func__)
\end{dhkcode}

예상 출력:

\begin{dhkzsh}
True
True
\end{dhkzsh}

이것은 \zsh{staticmethod}가 속성 조회 시 원본 함수 객체를 반환하고, 인스턴스나 클래스를 자동 인수로 바인딩하지 않기 때문이다.

따라서 다음 두 호출은 모두 정상적으로 실행된다.

\begin{dhkcode}
print(Calculator.add(10, 20))
print(obj.add(10, 20))
\end{dhkcode}

예상 출력:

\begin{dhkzsh}
30
30
\end{dhkzsh}

여기서 \zsh{obj.add(10, 20)}은 일반적인 인스턴스 메서드 호출과 다르다.

인스턴스 \zsh{obj}가 첫 번째 인수로 자동 전달되지 않는다.

결국 \zsh{staticmethod}의 정확한 의미는 다음과 같다.

\begin{dhksmry}[centered]
클래스 네임스페이스에 함수를 저장하되, 일반적인 인스턴스 

메서드 바인딩을 적용하지 않도록 하는 디스크립터다.
\end{dhksmry}

\subsubsection{\zsh{property} - 속성 조회가 함수 호출을 유발하는 구조}

이번에는 \zsh{property}를 살펴보자.

\begin{dhkcode}
class Account:

    def __init__(self, balance):
        self._balance = balance

    @property
    def balance(self):
        return self._balance


account = Account(100)
\end{dhkcode}

이 코드에서

\begin{dhkcode}
@property
def balance(self):
    return self._balance
\end{dhkcode}

는 개념적으로 다음과 대응한다.

\begin{dhkcode}
def balance(self):
    return self._balance

balance = property(balance)
\end{dhkcode}

따라서 클래스 네임스페이스에는 함수 객체가 아니라 \zsh{property} 객체가 저장된다.

\begin{dhkcode}
raw = Account.__dict__["balance"]

print(type(raw))
print(type(raw.fget))
\end{dhkcode}

예상 출력:

\begin{dhkzsh}
<class 'property'>
<class 'function'>
\end{dhkzsh}

\zsh{raw.fget}은 getter 함수 객체를 참조한다.

이제 다음 표현식을 평가한다고 하자.

\begin{dhkcode}
account.balance
\end{dhkcode}

일반적인 함수 디스크립터였다면 바운드 메서드가 반환되었겠지만, 

\zsh{property}는 다르게 동작한다.

\zsh{property.__get__}이 getter 함수를 호출하고, 그 함수의 반환값을 속성 조회 결과로 반환한다.

따라서 다음 결과를 얻는다.

\begin{dhkcode}
print(account.balance)
\end{dhkcode}

예상 출력:

\begin{dhkzsh}
100
\end{dhkzsh}

핵심 실행 흐름은 다음과 같다.

\begin{dhksmry}[centered]
account.balance
\da
Account.__dict__["balance"]
\da
property 객체의 __get__
\da
getter(account)
\da
account._balance 조회
\da
100 반환
\end{dhksmry}

\begin{dhknote}
\begin{dhksmry}[centered]
property는 데이터 디스크립터다
\end{dhksmry}
\end{dhknote}

속성 탐색 우선순위를 생각해보자.

일반적인 인스턴스 속성 조회에서는 데이터 디스크립터가 인스턴스의 

\zsh{__dict__}보다 우선한다.

\zsh{property}는 setter가 정의되지 않았더라도 데이터 디스크립터에 해당한다.

다음 코드로 확인할 수 있다.

\begin{dhkcode}
account.__dict__["balance"] = 999

print(account.__dict__["balance"])
print(account.balance)
\end{dhkcode}

예상 출력:

\begin{dhkzsh}
999
100
\end{dhkzsh}

인스턴스 딕셔너리에 \zsh{"balance"}라는 키가 존재하지만

\zsh{account.balance}를 조회하면 \zsh{property}가 우선 적용된다.

따라서 getter가 실행되어 \zsh{self._balance}의 값인 100이 반환되는 것이다.

\begin{dhknote}
\begin{dhksmry}[centered]
@balance.setter는 무엇을 반환할까?
\end{dhksmry}
\end{dhknote}

\begin{dhkcode}
class Account:

    def __init__(self, balance):
        self._balance = balance

    @property
    def balance(self):
        return self._balance

    @balance.setter
    def balance(self, value):
        if value < 0:
            raise ValueError("잔액은 음수일 수 없습니다.")

        self._balance = value
\end{dhkcode}

여기서 \zsh{@balance.setter}는 기존 \zsh{property} 객체에 setter를 제자리에서 추가하는 것과는 다르다.

\zsh{property.setter()}는 지정된 setter를 포함하는 새로운 \zsh{property} 객체를 반환한다.

그 반환 객체가 다시 클래스 네임스페이스의 \zsh{balance} 이름에 바인딩된다.

최종 객체를 확인해 보자.

\begin{dhkcode}
descriptor = Account.__dict__["balance"]

print(type(descriptor.fget))
print(type(descriptor.fset))
\end{dhkcode}

예상 출력:

\begin{dhkzsh}
<class 'function'>
<class 'function'>
\end{dhkzsh}

이제 다음 표현식은

\begin{dhkcode}
account.balance = 200
\end{dhkcode}

일반적인 인스턴스 딕셔너리 할당이 아니라 \zsh{property.__set__}을 통한 setter 호출로 처리된다.

따라서 \zsh{property}는 함수 호출을 속성 읽기·쓰기 연산에 연결하는 디스크립터라고 이해하면 된다.

\subsubsection{사용자 정의 데코레이터를 인스턴스 메서드에 적용하기}

이제 이번 학습의 핵심으로 들어가 본다.

일반적인 함수 데코레이터를 인스턴스 메서드에 적용하면 어떤 일이 발생할까?

\begin{dhkcode}
from functools import wraps


def trace(func):

    @wraps(func)
    def wrapper(*args, **kwargs):
        print("TRACE:", func.__name__)
        return func(*args, **kwargs)

    return wrapper


class Calculator:

    def __init__(self, base):
        self.base = base

    @trace
    def add(self, value):
        return self.base + value


obj = Calculator(10)

print(obj.add(5))
\end{dhkcode}

예상 출력:

\begin{dhkzsh}
TRACE: add
15
\end{dhkzsh}

이 코드는 세 단계로 분석해야 한다.

\begin{dhknote}
\begin{dhksmry}[centered]
1. 클래스 본문에서 데코레이터 적용
\end{dhksmry}
\end{dhknote}

먼저 원본 \zsh{add} 함수 객체가 생성된다.

그다음 \zsh{trace}에 원본 함수 객체가 전달된다.

\begin{dhkcode}
add = trace(add)
\end{dhkcode}

\zsh{trace}는 wrapper 함수 객체를 생성하고 반환한다.

따라서 클래스 네임스페이스에 저장되는 객체는 원본 \zsh{add}가 아니라 wrapper다.

\begin{dhknote}
\begin{dhksmry}[centered]
2. 인스턴스에서 add 조회
\end{dhksmry}
\end{dhknote}

\begin{dhkcode}
obj.add
\end{dhkcode}

클래스 네임스페이스에서 발견되는 객체는 wrapper 함수 객체다.

그런데 wrapper 역시 일반적인 Python 함수 객체이므로 \zsh{__get__}을 제공한다.

따라서 인스턴스 속성 조회에서 wrapper가 \zsh{obj}에 바인딩된다.

즉, \zsh{obj.add}는 wrapper를 감싸는 바운드 메서드 객체다.

\begin{dhknote}
\begin{dhksmry}[centered]
3. 바운드 메서드 호출
\end{dhksmry}
\end{dhknote}

\begin{dhkcode}
obj.add(5)
\end{dhkcode}

바운듣 메서드는 다음과 같은 호출을 수행한다.

\begin{dhkcode}
wrapper(obj, 5)
\end{dhkcode}

따라서 wrapper 내부에서 다음 바인딩이 형성된다.

\begin{dhkzsh}
args = (obj, 5)
kwargs = {}
\end{dhkzsh}

이후 wrapper는 클로저에 저장된 원본 \zsh{add} 함수 객체를 호출한다.

\begin{dhkcode}
func(*args, **kwargs)
\end{dhkcode}

결국 원본 함수는 다음 인수를 전달받아 실행된다.

\begin{dhkcode}
original_add(obj, 5)
\end{dhkcode}

원본 함수에서 \zsh{self}는 \zsh{obj}를 참조하고 \zsh{value}는 5를 참조한다.

따라서 \zsh{obj.base + value}의 결과인 15를 반환한다.

\begin{dhknote}
\begin{dhksmry}[centered]
바인딩 구조를 검증해 보자
\end{dhksmry}
\end{dhknote}

\begin{dhkcode}
bound = obj.add

wrapper_function = Calculator.__dict__["add"]
original_function = wrapper_function.__wrapped__

print(bound.__self__ is obj)
print(bound.__func__ is wrapper_function)
print(wrapper_function is original_function)
\end{dhkcode}

예상 출력:

\begin{dhkzsh}
True
True
False
\end{dhkzsh}

이 결과는 다음 사실을 보여준다.

\begin{dhksmry}[centered]
인스턴스에 바인딩되는 함수는 원본 add가 아니라

데코레이터가 반환한 wrapper 함수 객체다.
\end{dhksmry}

그리고 wrapper가 원본 함수에 \zsh{self}를 포함한 인수들을 전달하는 것이다.

\subsubsection{\zsh{@classmethod}와 사용자 정의 데코레이터의 적용 순서}

\zsh{A.10}에서 학습한 데코레이터 적용 순서는 클래스 메서드에서도 그대로 성립한다.

그런데 \zsh{classmethod}가 반환하는 객체는 일반 함수 객체가 아니라 디스크립터이기 때문에, 적용 순서를 바꾸면 결과가 크게 달라질 수 있다.

먼저 다음 코드를 보자.

\begin{dhkcode}
from functools import wraps


def trace(func):

    @wraps(func)
    def wrapper(*args, **kwargs):
        print("TRACE:", func.__name__)
        return func(*args, **kwargs)

    return wrapper


class Service:

    def __init__(self, value):
        self.value = value

    @classmethod
    @trace
    def create(cls, value):
        return cls(value)


service = Service.create(10)

print(type(service).__name__)
print(service.value)
\end{dhkcode}

예상 출력:

\begin{dhkzsh}
TRACE: create
Service
10
\end{dhkzsh}

여기서 데코레이터는 아래에서 위로 적용된다.

따라서 다음 코드와 대응한다.

\begin{dhkcode}
create = classmethod(trace(create))
\end{dhkcode}

이 과정을 객체 생성 순서로 분석해 보자.

\begin{dhknote}
\begin{dhksmry}[centered]
첫 번째 단계
\end{dhksmry}
\end{dhknote}

원본 \zsh{create} 함수 개체가 생성된다.

\begin{dhknote}
\begin{dhksmry}[centered]
두 번째 단계
\end{dhksmry}
\end{dhknote}

\zsh{trace}가 원본 함수 객체를 전달받아 wrapper 함수 객체를 생성한다.

\begin{dhknote}
\begin{dhksmry}[centered]
세 번째 단계
\end{dhksmry}
\end{dhknote}

\zsh{classmethod}가 wrapper 함수 객체를 감싸는 디스크립터를 생성한다.

\begin{dhknote}
\begin{dhksmry}[centered]
네 번째 단계
\end{dhksmry}
\end{dhknote}

클래스 네임스페이스의 \zsh{create} 이름에 \zsh{classmethod} 객체가 저장된다.

이제 다음 표현식이 평가된다.

\begin{dhkcode}
Service.create
\end{dhkcode}

\zsh{classmethod}의 \zsh{__get__}이 wrapper를 \zsh{Service} 클래스에 바인딩한 메서드를 반환한다.

따라서 다음 호출은

\begin{dhkcode}
Service.create(10)
\end{dhkcode}

개념적으로 다음과 같다.

\begin{dhkcode}
wrapper(Service, 10)
\end{dhkcode}

Wrapper는 클로저에 저장된 원본 함수 객체를 호출한다.

\begin{dhkcode}
func(Service, 10)
\end{dhkcode}

결국 원본 \zsh{create}의 \zsh{cls}는 \zsh{Service} 클래스 객체를 참조하게 된다.

\begin{dhknote}
\begin{dhksmry}[centered]
내부 구조 검증
\end{dhksmry}
\end{dhknote}

\begin{dhkcode}
descriptor = Service.__dict__["create"]

bound = Service.create

wrapper_function = descriptor.__func__
original_function = wrapper_function.__wrapped__

print(type(descriptor))

print(bound.__self__ is Service)
print(bound.__func__ is wrapper_function)

print(wrapper_function is original_function)
\end{dhkcode}

예상 출력:

\begin{dhkzsh}
<class 'classmethod'>
True
True
False
\end{dhkzsh}

이 관계를 그림으로 나타내면 다음과 같다.

\begin{dhkgraph}[uniform]
node n1 (1,1) :: Service.__dict__["create"] // classmethod 객체
node n2 (2,1) :: trace의 wrapper 함수 객체 // 클로저의 func -> 원본 create
node n3 (3,1) :: 원본 create 함수 객체 // cls(value)를 실행

edge n1 -> n2 :: __func__
edge n2 -> n3 :: __wrapped__

explain :: Service.create -> wrapper가 클래스에 바인딩된 메서드
\end{dhkgraph}

여기서 \zsh{classmethod}는 원본 함수를 직접 감싸는 것이 아니라, \zsh{trace}가 반환한 wrapper를 감싸고 있다는 점이 중요하다.

\subsubsection{\zsh{@trace @classmethod}로 순서를 바꾸면?}

이번에는 데코레이터의 순서를 반대로 작성해 보자.

\begin{dhkcode}
class Service:

    @trace
    @classmethod
    def create(cls, value):
        return cls(value)
\end{dhkcode}

이 코드는 다음과 대응한다.

\begin{dhkcode}
create = trace(classmethod(create))
\end{dhkcode}

앞의 예제와 결정적인 차이가 있다.

이번에는 \zsh{trace}에 전달되는 객체가 일반 함수가 아니라 \zsh{classmethod} 객체다.

\zsh{trace}는 이를 \zsh{func}라는 이름으로 참조하고 ㅈrapper를 생성한다.

최종적으로 클래스 네임스페이스에 저장되는 객체는 \zsh{classmethod}가 아니라 일반 wrapper 함수 객체다.

따라서 \zsh{Service.create}를 조회할 때 \zsh{classmethod}의 클래스 바인딩이 적용되지 않는다.

게다가 wrapper 내부에서 다음 코드를 실행하면

\begin{dhkcode}
return func(*args, **kwargs)
\end{dhkcode}

\zsh{func}가 참조하는 \zsh{classmethod} 객체를 직접 호출하려고 한다.

일반적인 \zsh{classmethod} 객체는 직접 호출할 수 없으므로 \zsh{TypeError}가 발생한다.

두 구조를 비교해 보자.

\begin{dhknote}
\begin{dhksmry}[centered]
클래스 바인딩 유지
\end{dhksmry}
\end{dhknote}

\begin{dhkcode}
@classmethod
@trace
def create(cls):
    ...
\end{dhkcode}

\begin{dhksmry}
    - 최종 저장 객체 : classmethod(wrapper)

    - 속성 조회 : wrapper를 클래스에 바인딩
\end{dhksmry}

\begin{dhknote}
\begin{dhksmry}[centered]
일반적인 trace에서는 문제 발생
\end{dhksmry}
\end{dhknote}

\begin{dhkcode}
@trace
@classmethod
def create(cls):
    ...
\end{dhkcode}

\begin{dhksmry}
    - 최종 저장 객체 : 일반 wrapper 함수

    - 속성 조회 : 클래스 메서드 바인딩이 유지되지 않음
\end{dhksmry}

따라서 일반적인 함수 ㅈrapper 데코레이터를 클래스 메서드에 적용할 때는 보통 다음 구조가 적절하다.

\begin{dhkcode}
@classmethod
@trace
def create(cls):
    ...
\end{dhkcode}

다만 모든 데코레이터가 반드시 이 순서를 요구하는 것은 아니다.

\zsh{classmethod} 디스크립터를 인식하고 별도로 처리하도록 설계한 데코레이터라면 다른 구조도 가능하다.

판단 기준은 데코레이터가 어떤 타입의 객체를 전달받고 어떤 타입의 객체를 반환하는가이다.

\subsubsection{\zsh{staticmethod}와 데코레이터의 적용 순서}

정적 메서드에도 같은 원리를 적용할 수 있다.

\begin{dhkcode}
class Calculator:

    @staticmethod
    @trace
    def multiply(a, b):
        return a * b


obj = Calculator()

print(Calculator.multiply(3, 4))
print(obj.multiply(3, 4))
\end{dhkcode}

예상 출력:

\begin{dhkzsh}
TRACE: multiply
12
TRACE: multiply
12
\end{dhkzsh}

이 구조는 다음과 대응한다.

\begin{dhkcode}
multiply = staticmethod(trace(multiply))
\end{dhkcode}

먼저 원본 함수에 \zsh{trace}가 적용되어 wrapper가 생성되고, 

그 wrapper를 \zsh{staticmethod}가 감싼다.

따라서 클래스나 인스턴스를 통해 속성을 조회해도 wrapper 함수 객체가 그대로 반환된다.

인스턴스가 첫 번째 인수로 자동 전달되지 않는 것이다.

\begin{dhknote}
\begin{dhksmry}[centered]
순서를 반대로 작성하면?
\end{dhksmry}
\end{dhknote}

\begin{dhkcode}
class Calculator:

    @trace
    @staticmethod
    def multiply(a, b):
        return a * b
\end{dhkcode}

이번에는 다음과 대응한다.

\begin{dhkcode}
multiply = trace(staticmethod(multiply))
\end{dhkcode}

\zsh{trace}가 \zsh{staticmethod} 객체를 전달받아 일반 함수 wrapper를 반환한다.

여기서 \zsh{classmethod}와는 차이가 있다.

Python 3.10부터 \zsh{staticmethod} 객체 자체도 직접 호출할 수 있다.

따라서 다음과 같은 클래스 접근 호출은 동작할 수 있다.

\begin{dhkcode}
Calculator.multiply(3, 4)
\end{dhkcode}

하지만 인스턴스 접근은 문제가 달라진다.

\begin{dhkcode}
obj = Calculator()

obj.multiply(3, 4)
\end{dhkcode}

최종 클래스 속성에 저장된 것은 일반 wrapper 함수 객체이므로 인스턴스 바인딩이 발생한다.

Wrapper는 \zsh{obj}까지 포함하여 원본 쪽으로 인수를 전달하려고 한다.

결국 원본 \zsh{multiply(a, b)}에 세 개의 위치 인수가 전달되어 \zsh{TypeError}가 발생할 수 있다.

즉, \zsh{staticmethod} 객체가 호출 가능하다는 사실과 정적 메서드의 바인딩 규칙이 보존된다는 사실은 별개다.

이 경우에도 일반적으로 다음 순서가 의도에 맞다.

\begin{dhkcode}
@staticmethod
@trace
def multiply(a, b):
    ...
\end{dhkcode}

\subsubsection{\zsh{property}와 사용자 정의 데코레이터}

\zsh{property}에도 함수 wrapper를 적용할 수 있다.

\begin{dhkcode}
from functools import wraps

def trace(func):

    @wraps(func)
    def wrapper(*args, **kwargs):
        print("TRACE:", func.__name__)
        return func(*args, **kwargs)

    return wrapper


class Account:

    def __init__(self, balance):
        self._balance = balance

    @property
    @trace
    def balance(self):
        return self._balance


account = Account(100)

print(account.balance)
\end{dhkcode}

예상 출력:

\begin{dhkzsh}
TRACE: balance
100
\end{dhkzsh}

이 구조는 다음과 대응한다.

\begin{dhkcode}
balance = property(trace(balance))
\end{dhkcode}

먼저 원본 getter 함수에 \zsh{trace}가 적용되어 wrapper를 생성한다.

그다음 \zsh{property}가 wrapper를 getter로 저장한다.

따라서 \zsh{account.balance}를 조회하면 다음 호출 경로가 형성된다.

\begin{dhksmry}[centered]
account.balance
\da
property.__get__
\da
wrapper(account)
\da
원본 getter(account)
\da
100 반환
\end{dhksmry}

여기서 wrapper는 인스턴스에 바인딩된 메서드로 조회되는 것이 아니다.

\zsh{property}가 getter로 저장된 wrapper 함수 객체를 인스턴스 인수와 함께 직접 호출하는 구조다.

반대로 다음과 같이 작성하면

\begin{dhkcode}
@trace
@property
def balance(self):
    ...
\end{dhkcode}

\zsh{trace}가 \zsh{property} 객체를 전달받게 된다.

일반적인 \zsh{trace}는 전달받은 객체를 함수처럼 호출하도록 구현되어 있으므로 정상적인 getter 동작을 보장하지 못한다.

더구나 최종적으로 클래스 네임스페이스에 일반 wrapper 함수가 저장되므로

\zsh{property}의 데이터 디스크립터 동작도 유지되지 않는다.

여기서도 데코레이터의 적용 순서가 최종 클래스 속성의 종류를 결정한다는 원리가 적용된다.

\subsubsection{호출 가능한 객체와 메서드 바인딩은 같은 개념일까?}

이번에는 한 단계 더 깊이 들어가 보자.

\zsh{A.9}에서 클래스 인스턴스를 데코레이터로 사용할 수 있다는 사실을 학습했다.

그렇다면 함수 wrapper 대신 호출 가능한 클래스 인스턴스를 반환하면 메서드 바인딩도 자동으로 이루어질까?

다음 코드를 보자.

\begin{dhkcode}
class Counter:

    def __init__(self, func):
        self.func = func
        self.calls = 0

    def __call__(self, *args, **kwargs):
        self.calls += 1
        return self.func(*args, **kwargs)


class Example:

    @Counter
    def calculate(self, x):
        return x * 2


obj = Example()
\end{dhkcode}

클래스 본문에서는 개념적으로 다음 변환이 수행된다.

\begin{dhkcode}
calculate = Counter(calculate)
\end{dhkcode}

따라서 클래스 네임스페이스에 저장되는 객체는 함수 객체가 아니라 \zsh{Counter} 인스턴스다.

\begin{dhkcode}
print(type(Example.__dict__["calculate"]))

print(
    obj.calculate
    is Example.__dict__["calculate"]
)
\end{dhkcode}

예상 출력:

\begin{dhkzsh}
<class '__main__.Counter'>
True
\end{dhkzsh}

두 번재 결과에 주목해 보자.

\zsh{Counter}는 \zsh{__call__}을 구현했으므로 호출 가능한 객체다.

하지만 \zsh{__get__}은 구현하지 않았다.

따라서 일반적인 함수 객체처럼 인스턴스 바인딩을 수행하지 않는다.

이제 다음 호출을 생각해 보자.

\begin{dhkcode}
obj.calculate(10)
\end{dhkcode}

\zsh{obj.calculate}는 클래스에 저장된 \zsh{Counter} 인스턴스 자체를 반환한다.

그 객체의 \zsh{__call__}이 실행되고 다음 코드에 도달한다.

\begin{dhkcode}
return self.func(*args, **kwargs)
\end{dhkcode}

여기서 \zsh{self.func}은 원본 \zsh{calculate} 함수 객체를 참조한다.

원본 함수는 두 개의 인수를 요구한다.

\begin{dhkcode}
def calculate(self, x):
    ...
\end{dhkcode}

하지만 전달된 인수는 10 하나뿐이다.

따라서 \zsh{TypeError}가 발생한다.

\begin{dhknote}
\begin{dhksmry}[centered]
왜 일반 함수 wrapper는 정상적으로 동작했을까?
\end{dhksmry}
\end{dhknote}

두 객체의 프로토콜이 다르기 때문이다.

\begin{dhktable}
|구분|	일반 함수 wrapper|	위 예제의 Counter 인스턴스|
|:---|:---|:---|
|호출 가능 여부|	가능|	가능|
|__call__|	제공|	제공|
|__get__|	제공|	제공하지 않음|
|인스턴스 접근 결과|	바운드 메서드|	Counter 객체 자체|
|인스턴스 자동 전달|	수행	|수행하지 않음|
\end{dhktable}

이것은 중요한 객체 모델의 구분이다.

호출 가능한 객체(Callable)라는 사실만으로 디스크립터 바인딩까지 제공되는 것은 아니다.

\zsh{__call__}은 호출 프로토콜에 관여하고, \zsh{__get__}은 속성 조회 및 바인딩 프로토콜에 관여한다.

\subsubsection{클래스 기반 데코레이터에 \zsh{__get__}구현하기}

그렇다면 앞의 \zsh{Counter}를 어떻게 수정해야 할까?

가장 직접적인 방법은 \zsh{__get__}을 구현하는 것이다.

\begin{dhkcode}
class Counter:

    def __init__(self, func):
        self.func = func
        self.calls = 0

    def __call__(self, *args, **kwargs):
        self.calls += 1
        return self.func(*args, **kwargs)

    def __get__(self, instance, owner=None):

        if instance is None:
            return self

        def bound(*args, **kwargs):
            return self(instance, *args, **kwargs)

        return bound


class Example:

    @Counter
    def calculate(self, x):
        return x * 2


a = Example()
b = Example()

print(a.calculate(10))
print(b.calculate(20))

print(Example.__dict__["calculate"].calls)
\end{dhkcode}

예상 출력:

\begin{dhkzsh}
20
40
2
\end{dhkzsh}

이번에는 왜 정상적으로 동작할까?

\zsh{a.calculate}를 조회하면 \zsh{Counter} 인스턴스의 \zsh{__get__}이 호출된다.

\begin{dhkcode}
counter.__get__(a, Example)
\end{dhkcode}

이 함수는 \zsh{bound}라는 새로운 함수 객체를 생성하고 반환한다.

\zsh{bound}는 클로저를 통해 \zsh{instance}와 \zsh{self}를 참조한다.

여기서 \zsh{instance}는 \zsh{a}를, \zsh{self}는 \zsh{Counter} 인스턴스를 참조한다.

따라서 다음 호출은

\begin{dhkcode}
a.calculate(10)
\end{dhkcode}

다음 순서로 실행된다.

\begin{dhksmry}[centered]
a.calculate(10)
\da
bound(10)
\da
counter(a, 10)
\da
Counter.__call__(counter, a, 10)
\da
원본 calculate(a, 10)
\da
20 반환
\end{dhksmry}

여기서 한 가지 주의해야 한다.

이 구현이 반환한 \zsh{bound}는 일반적인 \zsh{method} 타입의 바운드 메서드가 아니다.

인스턴스 참조를 클로저로 유지하는 일반 함수 객체다.

따라서 일반적인 바운드 메서드의 \zsh{__self__}, \zsh{__func__} 구조와는 완전히 같지는 않다.

\begin{dhknote}
\begin{dhksmry}[centered]
호출 횟수는 어디에 저장될까?
\end{dhksmry}
\end{dhknote}

\begin{dhkcode}
print(Example.__dict__["calculate"].calls)
\end{dhkcode}

결과는 2다.

\zsh{a}와 \zsh{b}는 서로 다른 인스턴스지만, 클래스 네임스페이스에 저장된 하나의 \zsh{Counter} 객체를 공유하기 때문이다.

따라서 이 예제에서 \zsh{calls}는 인스턴스별 호출 횟수가 아니라 데코레이터 객체의 누적 호출 횟수다.

이 차이는 실전에서 상태를 가진 클래스 기반 데코레이터를 설계할 때 중요하다.

\subsubsection{디스크립터와 인스턴스 네임스페이스의 탐색 우선순위}

이번에는 \zsh{A.11}의 내용을 기존 속성 탐색 규칙과 연결해 본다.

일반적인 \zsh{object.__getattribute__} 기반의 인스턴스 속성 탐색에서는 다음 우선순위가 적용된다.

\begin{dhktable}
|우선순위|	탐색 대상|
|:---|:---|
|1|	클래스의 MRO에서 찾은 데이터 디스크립터|
|2|	인스턴스의 __dict__|
|3|	클래스의 MRO에서 찾은 비데이터 디스크립터 또는 일반 속성|
|4|	일반 탐색 실패 시 적용될 수 있는 __getattr__|
\end{dhktable}

일반적인 함수 객체는 비데이터 디스크립터다.

따라서 인스턴스의 \zsh{__dict__}에 같은 이름의 속성이 있으면 인스턴스 속성이 우선할 수 있다.

\begin{dhkcode}
class Example:

    def method(self):
        return "class method"


obj = Example()

obj.method = lambda: "instance attribute"

print(obj.method())
print(Example.method(obj))
\end{dhkcode}

예상 출력:

\begin{dhkzsh}
instance attribute
class method
\end{dhkzsh}

왜 이런 결과가 발생할까?

\zsh{obj.method}를 조회할 때 인스턴스 딕셔너리에 \zsh{"method"}가 존재하므로 그 객체가 반환된다.

이때 저장된 \zsh{lambda} 함수는 인스턴스 속성으로 직접 조회되는 것이므로 클래스 함수 디스크립터의 바인딩이 적용되지 않는다.

반면 \zsh{Example.method}는 클래스 네임스페이스에 저장된 원본 함수 객체를 조회한다.

따라서 \zsh{Example.method(obj)}는 원본 함수를 명시적으로 호출하는 형태다.

이것은 일반적인 함수 객체가 비데이터 디스크립터라는 사실과 연결된다.

반면 \zsh{property}는 데이터 디스크립터이므로 인스턴스 딕셔너리보다 우선한다.

즉, 데코레이터가 반환하는 객체가 데이터 디스크립터인지, 비데이터 디스크립터인지, 일반 객체인지에 따라 속성 탐색 결과가 달라질 수 있다.

\subsubsection{\zsh{inspect.signature()}와 바운드 메서드}

\zsh{A.8}에서 학습한 함수 메타데이터를 메서드 바인딩과 연결해 보자.

\begin{dhkcode}
import inspect
from functools import wraps

def trace(func):

    @wraps(func)
    def wrapper(*args, **kwargs):
        return func(*args, **kwargs)

    return wrapper


class Calculator:

    @trace
    def add(self, x: int, y: int = 1) -> int:
        return x + y


obj = Calculator()

print(inspect.signature(Calculator.add))
print(inspect.signature(obj.add))
\end{dhkcode}

예상 출력:

\begin{dhkzsh}
(self, x: int, y: int = 1) -> int
(x: int, y: int = 1) -> int
\end{dhkzsh}

왜 결과가 다를까?

\zsh{Calculator.add}는 클래스에서 조회한 wrapper 함수 객체다.

\zsh{@wraps(func)}가 설정한 \zsh{__wrapped__}를 따라 원본 함수의 시그니처를 확인할 수 있다.

따라서 첫 번째 결과에는 \zsh{self}가 포함된다.

반면 \zsh{obj.add}는 인스턴스에 바인딩된 메서드 객체다.

\zsh{inspect.signature()}는 바운드 메서드의 바인딩 상태를 고려해 이미 바인딩된 첫 번째 인수를 호출자에게 요구되는 인수 목록에서 제외한다.

여기서 서로 다른 세 메커니즘이 결합된다.

\begin{dhktable}
|메커니즘	|역할|
|:---|:---|
|functools.wraps|	메타데이터와 __wrapped__ 참조 설정|
|함수 디스크립터 __get__|	인스턴스 바운드 메서드 생성|
|inspect.signature|	메타데이터와 바인딩 상태를 분석하여 시그니처 구성|
\end{dhktable}

이 셋을 같은 기능으로 이해하면 안 된다.

특히 \zsh{wraps}가 \zsh{self}를 자동으로 전달하는 것은 아니다.

실제 바인딩은 디스크립터 프로토콜이 담당한다.

\subsubsection{PyTorch \zsh{nn.Module}과의 연결}

이제 지금까지 학습한 객체 모델을 PyTorch에 연결해 보자.

\begin{dhkcode}
import torch
from torch import nn
from functools import wraps


def trace(func):

    @wraps(func)
    def wrapper(*args, **kwargs):
        print("FORWARD")
        return func(*args, **kwargs)

    return wrapper


class Model(nn.Module):

    def __init__(self):
        super().__init__()
        self.linear = nn.Linear(2, 1)

    @trace
    def forward(self, x):
        return self.linear(x)


model = Model()

x = torch.tensor([[1.0, 2.0]])

y = model(x)

print(y.shape)
\end{dhkcode}

예상 출력:

\begin{dhkzsh}
FORWARD
torch.Size([1, 1])
\end{dhkzsh}

이 코드는 두 가지 호출 메커니즘을 결합한다.

먼저 \zsh{Model} 클래스 네임스페이스에는 \zsh{trace}가 반환한 wrapper 함수 객체가 \zsh{forward}라는 이름으로 저장된다.

따라서 \zsh{model.forward}를 조회하면 wrapper가 \zsh{model} 인스턴스에 바인딩된다.

이것은 일반적인 Python 디스크립터 프로토콜의 동작이다.

하지만 다음 표현식은

\begin{dhkcode}
model(x)
\end{dhkcode}

단순히 \zsh{model.forward(x)}로 치환되는 것은 아니다.

\zsh{nn.Module}의 호출 메커니즘이 관여하며, 필요한 경우 등록된 훅(Hook) 등의 처리를 수행한다.

그 과정에서 \zsh{forward}가 호출되면 wrapper가 실행되고, 

wrapper는 원본 \zsh{forward}를 호출한다.

따라서 다음 네 단계를 구별해야 한다.

\begin{dhksmry}[centered]
model(x)
\da
nn.Module의 호출 메커니즘
\da
model.forward 속성 조회 및 호출
\da
trace의 wrapper 실행
\da
원본 forward 실행
\da
self.linear(x)
\da
Tensor 반환
\end{dhksmry}

이 예제에서 \zsh{@trace}는 \zsh{nn.Module.__call__}을 대체하는 것이 아니다.

클래스의 \zsh{forward} 속성에 저장되는 함수 객체를 wrapper로 변경하는 것이다.

그리고 PyTorch는 Python의 일반적인 객체 모델 위에 자체적인 모듈 호출 및 훅 처리 기능을 추가한 것이다.

\subsubsection{전체 개념 비교}

이번 학습에서 다룬 객체들을 하나의 표로 정리해 보자.

\begin{dhktable}
|객체|	클래스 속성에 저장된 형태|	속성 조회 시 동작|	핵심 메커니즘|
|:---|:---|:---|:---|
|일반 함수|	function|	인스턴스 접근 시 바운드 메서드 생성|	함수의 __get__|
|classmethod|	classmethod 객체	|클래스 바운드 메서드 생성|	classmethod.//*__get__|
|staticmethod|	staticmethod 객체|	감싼 함수 객체 반환|	staticmethod.//*__get__|
|property|	property 객체|	getter 실행 또는 setter 처리|	property.//*__get__, __set__|
|일반 함수 wrapper|	function|	인스턴스 접근 시 wrapper 바인딩|	함수의 __get__|
|__call__만 가진 데코레이터 인스턴스|	사용자 정의 객체|	자동 메서드 바인딩 없음|	호출 프로토콜만 제공|
|__get__도 가진 데코레이터 인스턴스|	사용자 정의 디스크립터|	구현에 따라 바인딩 수행|	사용자 정의 __get__|
\end{dhktable}

이 표에서 가장 중요한 구분은 호출 가능 여부와 디스크립터 여부가 서로 다른 성질이라는 점이다.

어떤 객체가 호출 가능하더라도 메서드 바인딩을 수행하지 않을 수 있다.

반대로 \zsh{property}처럼 일반적인 함수 호출 방식으로 직접 사용하지 않는 객체도 속성 조회에 개입할 수 있다.

\subsubsection{핵심 정리와 판단 기준}

이번 학습의 전체 실행 구조를 정리해 보자.

\begin{dhksmry}
클래스 메서드에 데코레이터가 적용되는 전체 흐름

1. 클래스 본문 실행
   def 문장이 원본 함수 객체를 생성한다.

2. 데코레이터 적용
   데코레이터가 원본 함수 또는 다른 데코레이터의 반환 객체를 전달받아
   새로운 객체를 반환할 수 있다.

3. 클래스 네임스페이스 바인딩
   최종 반환 객체가 클래스 속성 이름에 바인딩된다.

4. 클래스 또는 인스턴스 속성 조회
   속성 탐색 규칙에 따라 객체를 찾고, 디스크립터라면 __get__ 등의
   동작이 적용될 수 있다.

5. 호출 또는 속성값 사용
   반환된 객체가 바운드 메서드인지, 일반 함수인지, 계산된 속성값인지에
   따라 이후 동작이 달라진다.
\end{dhksmry}

실제 코드를 분석할 때는 다음 다섯 가지를 판단하면 된다.

\begin{dhksmry}
1. 클래스 네임스페이스에 저장된 객체 타입이 무엇인가?
   일반 함수인지, wrapper인지, classmethod인지, property인지에
   따라 속성 조회가 달라진다.

2. 그 객체는 어떤 디스크립터 메서드를 제공하는가?
   __get__, __set__, __delete__의 존재 여부가 속성 탐색과
   바인딩에 영향을 준다.

3. 바운드 메서드가 반환되었다면 __self__와 __func__는 무엇을 참조하는가?
   이를 확인하면 어떤 객체가 자동 인수로 전달되는지 추적할 수 있다.

4. wrapper는 클로저를 통해 어떤 객체를 참조하는가?
   원본 함수인지, 다른 wrapper인지, 디스크립터 객체인지 구별해야 한다.

5. 데코레이터 적용 순서를 바꾸면 클래스에 저장되는 객체도 달라지는가?
   객체가 달라진다면 이후 속성 조회와 호출 규칙도 달라질 수 있다.
\end{dhksmry}

이번 학습의 핵심 결론은 다음과 같다.

\begin{dhksmry}[centered]
데코레이터는 함수 정의 과정에서 객체를 변환하고,

디스크립터는 이후 속성 조회 과정에서 그 객체의 접근 결과를 결정한다.
\end{dhksmry}

이 두 메커니즘을 분리하면 인스턴스 메서드, 클래스 메서드, 정적 메서드, 프로퍼티, 사용자 정의 클래스 기반 데코레이터를 모두 같은 Python 객체 모델로 설명할 수 있다.
